\begin{proof}[Proof of \cref{obj:lem:universal-heavy-cell-rate}]
Write
\[
L=\log(en),\qquad m_j=|I_j|\quad (j=0,1),
\]
and let \(N_0\) denote the calibration cutoff in \cref{obj:def:hybrid-estimator-handle}. For
\(n\ge N_0\), the implemented heavy set \(\widehat{\mathcal H}_n\) is the pilot set at
threshold \(256\):
\[
\widehat{\mathcal H}
=
\{k:\lfloor 256L\rfloor<N_k^{(0)}\}.
\]
The corresponding pilot-bad event is
\[
\mathcal B
=
\left\{
\neg\left(\forall k\in\widehat{\mathcal H},\,\frac{256L}{2m_0}\le p_k\right)
\mathrel{\text{or}}
\neg\left(\forall k\notin\widehat{\mathcal H},\,p_k\le \frac{2\cdot256L}{m_0}\right)
\right\}.
\]

For every sample, the empirical arm ratios in \(\widehat h_k\) use totalized division, so the
ratio is \(0\) when the corresponding arm count is zero; with this convention each such ratio
lies in \([0,1]\). Hence the absolute empirical ratio contribution of category \(k\) is at most
\(N_k^{(1)}/m_1\), and summing over the selected categories gives
\[
|\widehat T_{\mathcal H}|
\le
\sum_{k\in\widehat{\mathcal H}_n}\frac{N_k^{(1)}}{m_1}
\le
\sum_k\frac{N_k^{(1)}}{m_1}
=
\frac{m_1}{m_1}
\le1.
\]
For an overlap law \(P\), the category vector \(q_k\) from \cref{obj:def:ate-functional}
lies in the overlap cone of \cref{obj:def:overlap-cone}. If \(p_k=0\), then \(q_k=0\) and
\cref{obj:def:cell-functional} sets \(\phi(q_k)=0\). If \(p_k>0\), the arm masses in
\(q_k\) are positive and the displayed formula in \cref{obj:def:cell-functional} gives
\[
\phi(q_k)=p_k(\mu_{1k}-\mu_{0k}),
\]
where the binary-outcome regressions \(\mu_{0k}\) and \(\mu_{1k}\) lie in \([0,1]\). Thus, in
both cases, \(|\phi(q_k)|\le p_k\). Since the category masses sum to one,
\[
\left|\sum_{k\in\widehat{\mathcal H}_n}\phi(q_k)\right|
\le
\sum_{k\in\widehat{\mathcal H}_n}p_k
\le
\sum_k p_k
=1.
\]
Thus the squared heavy-cell error is bounded by \(4\) pointwise for the overlap laws considered
below. The fixed-threshold pilot sandwich from \cref{obj:lem:pilot-sandwich}, applied with
dimension constant \(1\), gives constants \(C_0>0\) and \(N_{\mathrm{bad}}\) such that, whenever
\(n\ge N_{\mathrm{bad}}\) and \(d\le nL\),
\[
P^{\otimes n}(\mathcal B)\le C_0 n^{-4}.
\]
With \(C_{\mathrm{bad}}=4C_0\), the range bound gives
\[
\int_{\mathcal B}
\left(
\widehat T_{\mathcal H}
-
\sum_{k\in\widehat{\mathcal H}_n}\phi(q_k)
\right)^2
\,dP^{\otimes n}
\le C_{\mathrm{bad}}\,n^{-4}.
\] % lean: heavy_bad_pilot_rate

Set
\[
K^{\mathcal H}_\epsilon=\frac{16}{\epsilon}+12+\frac1{\epsilon^4},
\qquad
C_\epsilon=4K^{\mathcal H}_\epsilon+C_{\mathrm{bad}},
\qquad
\rho_\epsilon=1,
\]
and take
\[
N_\epsilon=\max\{8,N_0,N_{\mathrm{bad}}\}.
\] % lean: universal_heavy_cell_rate
These constants satisfy the required positivity conditions.

Assume \(n\ge N_\epsilon\), \(d\le \rho_\epsilon n\log n\), and
\(\mu_n\in\mathcal E_{n,d,\epsilon}\). The experiment-class assumption gives
\(\mu_n=P^{\otimes n}\) and \(\epsilon\)-overlap for \(P\). Also
\(n\ge8\), \(n\ge N_0\), \(n\ge N_{\mathrm{bad}}\), and, since \(\log n\le L\),
\[
d\le nL.
\]
Let
\[
B_- = \frac{256L}{2m_0}.
\]
On \(\mathcal B^c\), the first half of the pilot sandwich and the equality
\(\widehat{\mathcal H}_n=\widehat{\mathcal H}\) give
\[
k\in\widehat{\mathcal H}_n
\quad\Longrightarrow\quad
B_-\le p_k .
\] % lean: heavy_cell_mass_lower_of_good_pilot

Fix a deterministic category set \(H\) satisfying \(p_k\ge B_-\) for every \(k\in H\), and fix
\(a\in\{0,1\}\). Define
\[
S_a(H)
=
\sum_{k\in H}
\frac{N_k^{(1)}}{m_1}
\frac{N_{ak1}^{(1)}}{N_{ak}^{(1)}}
\mathbf 1\{N_{ak}^{(1)}>0\},
\qquad
\theta_a(H)=\sum_{k\in H}p_k\mu_{ak},
\]
where the displayed ratio is taken to be \(0\) when \(N_{ak}^{(1)}=0\). The empirical
category-mass centering is
\[
\sum_{k\in H}\frac{N_k^{(1)}}{m_1}\mu_{ak}.
\]
The pointwise decomposition about this centering is
\[
S_a(H)-\sum_{k\in H}\frac{N_k^{(1)}}{m_1}\mu_{ak}
=
\frac1{m_1}
\left[
\sum_{k\in H}
N_k^{(1)}\frac{\mathbf 1\{N_{ak}^{(1)}>0\}}{N_{ak}^{(1)}}
\{N_{ak1}^{(1)}-\mu_{ak}N_{ak}^{(1)}\}
-
\sum_{k\in H}\mu_{ak}N_k^{(1)}\mathbf 1\{N_{ak}^{(1)}=0\}
\right].
\] % lean: fixedHeavyArm_noise_decomposition

The first bracket is the aggregate centered ratio residual \(\mathcal R_a(H)\) of
\cref{obj:lem:heavy-residual-bound}, formed on the estimation split: the index set is \(I_1\),
so \(m=m_1\); the category set \(H\) is deterministic and carries \(p_k\ge B_->0\); and \(P\)
has overlap at level \(\epsilon\). \Cref{obj:lem:heavy-residual-bound} therefore gives
\[
\int
\left[
\sum_{k\in H}
N_k^{(1)}\frac{\mathbf 1\{N_{ak}^{(1)}>0\}}{N_{ak}^{(1)}}
\{N_{ak1}^{(1)}-\mu_{ak}N_{ak}^{(1)}\}
\right]^2
\,dP^{\otimes n}
\le
\frac{2m_1}{\epsilon}\sum_{k\in H}p_k .
\] % lean: integral_fixedRatioResidual_sq_le

For the second bracket, \(0\le\mu_{ak}\le1\), so its square is bounded by the square of the
missing-arm count
\[
\sum_{k\in H}N_k^{(1)}\mathbf 1\{N_{ak}^{(1)}=0\}.
\]
Since \(n\ge8\) we have \(m_1\ge3\), and the categories in \(H\) have mass at least \(B_->0\), so
that count meets the hypotheses of \cref{obj:lem:missing-arm-envelope} on the estimation split,
with \(m=m_1\) and mass floor \(B=B_-\). Combining the domination just recorded with the second
moment of \cref{obj:lem:missing-arm-envelope} gives
\[
\int
\left[
\sum_{k\in H}\mu_{ak}N_k^{(1)}\mathbf 1\{N_{ak}^{(1)}=0\}
\right]^2
\,dP^{\otimes n}
\le
m_1+m_1^2
\left[
\frac{|H|}
{\{((m_1-2)/2)\epsilon\}^2 B_-}
\right]^2 .
\] % lean: fixedMissingOutcomeBias_sq_le integral_fixedMissingCount_sq_le
Together with \((x-y)^2\le2x^2+2y^2\), the two bracket bounds give
\[
\int
\left(
S_a(H)-\sum_{k\in H}\frac{N_k^{(1)}}{m_1}\mu_{ak}
\right)^2
\,dP^{\otimes n}
\le
\frac{4\sum_{k\in H}p_k}{m_1\epsilon}
+
\frac2{m_1}
+
2\left[
\frac{|H|}
{\{((m_1-2)/2)\epsilon\}^2 B_-}
\right]^2 .
\] % lean: integral_fixedHeavyArm_noise_sq_le

For the empirical category-mass fluctuation, use the one-observation score
\[
g_{a,H}(O)=\mu_{aX}\mathbf 1\{X\in H\}.
\]
Its second-split average is the empirical category-mass centering, its expectation is
\(\theta_a(H)\), and \(0\le g_{a,H}\le1\). The iid bounded-score variance bound therefore gives
\[
\int
\left(
\sum_{k\in H}\frac{N_k^{(1)}}{m_1}\mu_{ak}-\theta_a(H)
\right)^2
\,dP^{\otimes n}
\le \frac1{m_1}.
\] % lean: integral_fixedEmpiricalMassArm_sub_target_sq_le
Splitting \(S_a(H)-\theta_a(H)\) at the empirical category-mass centering yields
\[
\int \{S_a(H)-\theta_a(H)\}^2\,dP^{\otimes n}
\le
\frac{8\sum_{k\in H}p_k}{m_1\epsilon}
+
\frac6{m_1}
+
4\left[
\frac{|H|}
{\{((m_1-2)/2)\epsilon\}^2 B_-}
\right]^2 .
\] % lean: integral_fixedHeavyArm_error_sq_le

The probability-mass identities give
\[
\sum_{k\in H}p_k\le\sum_k p_k=1,
\]
and the finite-set cardinality bound gives \(|H|\le d\). The split-size calibration for
\(n\ge8\) gives
\[
\frac n2\le m_1,
\qquad
\frac n8\le\frac{m_1-2}{2},
\qquad
\frac{256L}{n}\le B_-,
\]
and \(\log n\le L\). Hence
\[
\frac{8\sum_{k\in H}p_k}{m_1\epsilon}
+
\frac6{m_1}
+
4\left[
\frac{|H|}
{\{((m_1-2)/2)\epsilon\}^2 B_-}
\right]^2
\le
K^{\mathcal H}_\epsilon
\left\{
\frac1n+\frac{d^2}{n^2(\log n)^2}
\right\}.
\] % lean: heavy_fixed_envelope_rate

Let \(V=K^{\mathcal H}_\epsilon\{n^{-1}+d^2/[n^2(\log n)^2]\}\). The preceding deterministic bound applies
uniformly to every fixed \(H\) whose elements have mass at least \(B_-\). To pass to the random
pilot-heavy set on \(\mathcal B^c\), partition the pilot sample space into the fibers on which
\(\widehat{\mathcal H}_n=H\). The partition is carried out on the ambient iid sequence, whose
first \(n\) coordinates have law \(P^{\otimes n}\) by \cref{obj:lem:ambient-sample-marginal}, so
every integral below may equally be read under either law. On each contributing fiber, the
good-pilot implication gives
\(p_k\ge B_-\) for every \(k\in H\). The fixed-\(H\) arm error is a function of the estimation
split, while the fiber is determined by the pilot split, so the product law factors the restricted
integral into the fiber probability times the fixed-\(H\) product-law integral. Summing over the
eligible fibers and using that their probabilities sum to at most one gives, for each
\(a\in\{0,1\}\),
\[
\int_{\mathcal B^c}
\{S_a(\widehat{\mathcal H}_n)-\theta_a(\widehat{\mathcal H}_n)\}^2
\,dP^{\otimes n}
\le V.
\] % lean: heavy_good_arm_setIntegral_le

The heavy contribution and target decompose into their two arm components:
\[
\widehat T_{\mathcal H}
-
\sum_{k\in\widehat{\mathcal H}_n}\phi(q_k)
=
\{S_1(\widehat{\mathcal H}_n)-\theta_1(\widehat{\mathcal H}_n)\}
-
\{S_0(\widehat{\mathcal H}_n)-\theta_0(\widehat{\mathcal H}_n)\}.
\] % lean: fixedHeavyContribution_eq_arm_sub fixedTargetHeavy_eq_arm_sub
Thus \((x-y)^2\le2x^2+2y^2\) gives
\[
\int_{\mathcal B^c}
\left(
\widehat T_{\mathcal H}
-
\sum_{k\in\widehat{\mathcal H}_n}\phi(q_k)
\right)^2
\,dP^{\otimes n}
\le
4K^{\mathcal H}_\epsilon
\left\{
\frac1n+\frac{d^2}{n^2(\log n)^2}
\right\}.
\] % lean: fixedHeavy_error_sq_le_two_arms

The bad-event contribution satisfies
\[
\int_{\mathcal B}
\left(
\widehat T_{\mathcal H}
-
\sum_{k\in\widehat{\mathcal H}_n}\phi(q_k)
\right)^2
\,dP^{\otimes n}
\le
C_{\mathrm{bad}} n^{-4}
\le
C_{\mathrm{bad}}
\left\{
\frac1n+\frac{d^2}{n^2(\log n)^2}
\right\},
\] % lean: heavy_bad_pilot_rate
because \(n^{-4}\le n^{-1}\) for \(n\ge8\) and \(n^{-1}\le n^{-1}+d^2/[n^2(\log n)^2]\). Adding the
integrals over \(\mathcal B\) and \(\mathcal B^c\), and using \(\mu_n=P^{\otimes n}\), yields
\[
\int
\left(
\widehat T_{\mathcal H}
-
\sum_{k\in\widehat{\mathcal H}_n}\phi(q_k)
\right)^2
\,d\mu_n
\le
(4K^{\mathcal H}_\epsilon+C_{\mathrm{bad}})
\left\{
\frac1n+\frac{d^2}{n^2(\log n)^2}
\right\},
\]
which is the claimed bound with the chosen \(C_\epsilon\). % lean: universal_heavy_cell_rate
\end{proof}
